% ============================ CHAPITRE 2 =============================
\chapter{Analyse, conception et architecture de la solution}
\label{chap:conception}

\section{Introduction}

Ce chapitre traduit le positionnement établi au chapitre précédent en une conception
logicielle complète. Après une présentation générale du système et la formalisation des
besoins fonctionnels et non fonctionnels, il décrit le pipeline de traitement, l'architecture
globale à trois niveaux (frontend React, backend FastAPI, couche de persistance), le détail
des composants du backend, les flux entre composants, la modélisation UML, le modèle de
données, et s'achève sur la justification argumentée de chaque choix technologique.

\section{Présentation générale du projet}

Le système développé, baptisé \emph{Assistant RAG Multimodal}, est une application web
complète permettant à un utilisateur authentifié de :

\begin{itemize}
  \item téléverser des documents techniques (PDF, Word, Excel) et les organiser en
        \emph{groupes documentaires} (par exemple « Research Paper », « Technical Manual ») ;
  \item suivre en temps réel les étapes d'ingestion de chaque document (analyse, découpage,
        description des figures, vectorisation, indexation) ;
  \item consulter chaque document indexé : rendu Markdown du contenu extrait, galerie des
        figures avec leurs descriptions générées ;
  \item poser des questions en langage naturel dans une interface conversationnelle
        multi-conversations, avec choix du périmètre de recherche (tout le corpus, sélection
        manuelle de documents, ou routage automatique par LLM) ;
  \item obtenir des réponses \emph{sourcées} : chaque affirmation cite ses sources
        \texttt{[Source N]}, et lorsque la source est une figure, l'image elle-même est
        affichée sous la réponse.
\end{itemize}

Le corpus de démonstration comprend des articles de recherche fondateurs en apprentissage
profond (\emph{Attention Is All You Need}, BERT, ResNet, ViT, YOLO) et des manuels
d'ingénierie électrique (symboles de schémas, lecture de plans), deux familles choisies
pour leur richesse en figures techniques.

\section{Besoins fonctionnels}
\label{sec:besoins-fonctionnels}

\begin{table}[!htb]
\centering\small
\caption{Synthèse des besoins fonctionnels.}
\label{tab:bf}
\begin{tabularx}{\textwidth}{@{}lX@{}}
\toprule
\textbf{Réf.} & \textbf{Besoin} \\
\midrule
BF-01 & Téléverser des documents PDF, DOCX et XLSX, avec affectation à un groupe documentaire \\
BF-02 & Ingestion automatique : analyse structurelle, découpage sémantique, extraction et description des figures, vectorisation, indexation \\
BF-03 & Suivi de l'état d'ingestion (étape courante, succès, erreur) document par document \\
BF-04 & Gestion des documents : consultation détaillée (Markdown + figures), renommage, changement de groupe, duplication, suppression, téléchargement du fichier source \\
BF-05 & Gestion des groupes documentaires : création, renommage, suppression (avec repli des documents vers le groupe par défaut) \\
BF-06 & Poser une question en langage naturel et recevoir une réponse générée, ancrée dans le corpus, avec citations \texttt{[Source N]} \\
BF-07 & Restitution multimodale : affichage des figures sources sous la réponse lorsqu'elles sont citées \\
BF-08 & Périmètre de recherche : corpus entier, documents sélectionnés manuellement, ou routage automatique de la question vers les groupes pertinents par LLM \\
BF-09 & Historique de conversations par utilisateur : création, consultation, reprise, suppression \\
BF-10 & Authentification des utilisateurs (inscription, connexion) et cloisonnement des conversations par utilisateur \\
\bottomrule
\end{tabularx}
\end{table}

\section{Besoins non fonctionnels}
\label{sec:besoins-non-fonctionnels}

\begin{table}[!htb]
\centering\small
\caption{Synthèse des besoins non fonctionnels.}
\label{tab:bnf}
\begin{tabularx}{\textwidth}{@{}lX@{}}
\toprule
\textbf{Réf.} & \textbf{Besoin} \\
\midrule
BNF-01 & \textbf{Fiabilité des réponses} : le LLM ne doit répondre qu'à partir des sources récupérées et déclarer explicitement l'absence de réponse plutôt que d'inventer \\
BNF-02 & \textbf{Traçabilité} : toute réponse doit être vérifiable (document, page, extrait ou figure) \\
BNF-03 & \textbf{Compatibilité environnement contraint} : fonctionnement sans droits administrateur, sans service système à installer, à travers un proxy d'inspection TLS \\
BNF-04 & \textbf{Robustesse de l'ingestion} : un document en erreur ne doit ni bloquer la file ni corrompre l'index ; les erreurs sont enregistrées et affichées \\
BNF-05 & \textbf{Stabilité sous charge locale} : les téléversements en rafale ne doivent pas saturer la machine ni provoquer de contention d'écriture (sérialisation des ingestions) \\
BNF-06 & \textbf{Sécurité} : vérification cryptographique des jetons de session (RS256/JWKS), contrôle de propriété des conversations, CORS restrictif, validation des types de fichiers \\
BNF-07 & \textbf{Portabilité cloud} : déploiement cible sur OCI sans refonte (stockage embarqué, cache commutable vers un Redis managé par simple variable d'environnement) \\
BNF-08 & \textbf{Qualité mesurable} : le système doit être évaluable quantitativement (métriques RAGAS reproductibles) \\
BNF-09 & \textbf{Coût maîtrisé} : recours aux offres gratuites/économiques des fournisseurs de modèles (Mistral, DeepInfra via Hugging Face, Groq pour l'évaluation) \\
\bottomrule
\end{tabularx}
\end{table}

\section{Principe général du pipeline}
\label{sec:pipeline}

Le système s'organise en deux chaînes asynchrones l'une de l'autre : une chaîne
d'\textbf{ingestion} (hors ligne, déclenchée par chaque téléversement) et une chaîne de
\textbf{restitution} (en ligne, déclenchée par chaque question). La
figure~\ref{fig:pipeline-global} en donne la vue d'ensemble.

\rapportfig[0.98]{pipeline_global.pdf}{Pipeline fonctionnel complet : chaîne d'ingestion
(à gauche) et chaîne de restitution (à droite), articulées autour de l'index unifié
ChromaDB + BM25.}{fig:pipeline-global}

\subsection{Chaîne d'ingestion}
\begin{enumerate}
  \item \textbf{Analyse structurelle (Docling)} : le document est converti en une
        représentation structurée (sections, paragraphes, tableaux, figures avec leurs
        légendes et numéros de pages). Le backend d'analyse PDF retenu est
        \texttt{pypdfium2}, plus robuste en mémoire que le backend par défaut ;
  \item \textbf{Export Markdown} : le contenu structuré est exporté en Markdown et conservé
        pour la consultation dans l'interface ;
  \item \textbf{Découpage sémantique} : le \emph{HybridChunker} de Docling produit des
        chunks respectant la structure logique et la limite de 512 tokens, chacun
        \emph{contextualisé} (préfixé par ses titres de sections) et rattaché à sa page ;
  \item \textbf{Extraction et description des figures} : chaque image est sauvegardée sur
        disque puis décrite par le VLM Gemma-3 27B ; la légende d'origine trouvée par
        Docling (« Figure 1 : ... ») est préfixée à la description générée, pour que la
        recherche lexicale sur « Figure 1 » ait une prise ;
  \item \textbf{Vectorisation} : chunks de texte et descriptions de figures sont plongés
        dans le même espace vectoriel par \texttt{mistral-embed} (1024 dimensions,
        vecteurs normalisés L2), par lots de 32 ;
  \item \textbf{Indexation} : les vecteurs, textes et métadonnées (document, page, type
        \texttt{text}/\texttt{image}, chemin de l'image) sont insérés dans la collection
        ChromaDB ; l'index BM25 en mémoire est marqué à reconstruire.
\end{enumerate}

\subsection{Chaîne de restitution}
\begin{enumerate}
  \item \textbf{Périmètre} : si l'utilisateur n'a pas sélectionné de documents et a activé
        le périmètre automatique, un appel LLM route la question vers les groupes
        documentaires plausiblement pertinents (avec repli conservateur sur tout le corpus) ;
  \item \textbf{Recherche hybride} : la question est vectorisée, puis deux recherches sont
        lancées avec sur-échantillonnage (\(3 \times k\)) : dense (ChromaDB) et lexicale
        (BM25) ; leurs classements sont fusionnés par RRF et les \(k=5\) meilleurs chunks
        sont retenus ;
  \item \textbf{Génération contrainte} : les chunks numérotés (avec document et page) sont
        injectés dans le prompt de \texttt{mistral-small-latest}, dont les instructions
        imposent de ne répondre qu'à partir des sources, de citer \texttt{[Source N]} et de
        signaler l'absence de réponse ;
  \item \textbf{Restitution sourcée} : la réponse est renvoyée avec la liste structurée des
        sources ; celles de type figure portent l'URL de l'image, affichée par le frontend.
\end{enumerate}

\section{Architecture globale du système}
\label{sec:archi-globale}

L'architecture (figure~\ref{fig:archi-globale}) suit un découpage classique en trois
niveaux, enrichi d'une couche de services d'IA externes :

\begin{itemize}
  \item \textbf{Frontend} : SPA React 18 + TypeScript servie par Vite, intégrant le SDK
        Clerk pour l'authentification ; elle dialogue avec le backend en JSON/HTTP en
        joignant le jeton de session à chaque requête ;
  \item \textbf{Backend} : application FastAPI exposant trois familles de routes
        (\texttt{/api/documents}, \texttt{/api/groups}, \texttt{/api/chat} +
        \texttt{/api/conversations}), un serveur de fichiers statiques pour les images
        extraites, et un \emph{worker} d'ingestion sur file en arrière-plan ;
  \item \textbf{Persistance} : SQLite (métadonnées des documents, groupes, conversations,
        messages), ChromaDB persistant (vecteurs + textes + métadonnées), système de
        fichiers (fichiers sources, exports Markdown, images extraites), cache Redis
        commutable (FakeRedis en développement, OCI Cache en production) ;
  \item \textbf{Services d'IA externes} : API Mistral (embeddings \texttt{mistral-embed}
        et génération \texttt{mistral-small-latest}), API Hugging Face Inference /
        DeepInfra (captioning Gemma-3 27B), Clerk (JWKS de vérification des jetons),
        Groq (juge Llama 3.3 70B, pour l'évaluation uniquement).
\end{itemize}

\rapportfig[0.98]{architecture_globale.png}{Architecture globale : SPA React, API FastAPI,
couche de persistance embarquée et services d'IA externes.}{fig:archi-globale}

\section{Architecture du backend}
\label{sec:archi-backend}

Le backend suit une organisation en couches nettes (routes, services, stockage), sans
framework d'orchestration intermédiaire (figure~\ref{fig:classes-backend}) :

\begin{description}
  \item[\texttt{api/}] les routeurs FastAPI (\texttt{routes\_documents},
    \texttt{routes\_groups}, \texttt{routes\_chat}) : validation des entrées (schémas
    Pydantic), contrôle d'accès, orchestration des services ; aucune logique métier ;
  \item[\texttt{services/}] le cœur métier, un module par responsabilité :
    \texttt{docling\_parser} (analyse, découpage, extraction d'images),
    \texttt{ingestion} et \texttt{ingestion\_queue} (pipeline et sérialisation),
    \texttt{hf\_client} (embeddings multi-backends et captioning VLM),
    \texttt{mistral\_client} (embeddings, génération, routage de groupes),
    \texttt{vector\_store} (ChromaDB, BM25, fusion RRF),
    \texttt{clerk\_auth} (vérification JWT),
    \texttt{cache} (cache JSON fail-open),
    et les dépôts \texttt{documents\_repo}, \texttt{groups\_repo},
    \texttt{conversations\_repo} (accès SQLite) ;
  \item[\texttt{storage/}] la résolution centralisée des chemins de données
    (\texttt{uploads/}, \texttt{markdown/}, \texttt{images/}, \texttt{chroma/},
    \texttt{app.db}) ;
  \item[\texttt{config.py}] la configuration typée \texttt{pydantic-settings}, chargée
    depuis \texttt{.env} : clés d'API, choix du backend d'embeddings, chemins des modèles
    hors ligne, URL Redis, origines CORS.
\end{description}

\rapportfig[0.95]{uml_paquetages_backend.pdf}{Diagramme de paquetages du backend :
dépendances entre routes, services et stockage.}{fig:classes-backend}

\begin{keybox}[Un choix structurant : la file d'ingestion à worker unique]
L'ingestion (Docling + VLM + embeddings) est lourde en CPU et en mémoire. Le traitement
concurrent des téléversements en rafale saturait la machine et provoquait des contentions
d'écriture SQLite/ChromaDB allant jusqu'à l'échec de documents. L'architecture retenue,
une file en mémoire consommée par un \emph{unique} thread worker démarré avec
l'application, garantit un traitement strictement séquentiel, dans l'ordre d'arrivée,
quelle que soit la rafale. L'API reste, elle, parfaitement réactive : le téléversement
répond immédiatement et l'interface suit la progression par interrogation périodique.
\end{keybox}

\section{Architecture du frontend}

Le frontend est structuré par domaines fonctionnels : \texttt{components/auth} (pages et
garde d'authentification Clerk), \texttt{components/documents} (vue corpus, carte et détail
de document, dropzone de téléversement, galerie de figures), \texttt{components/chat}
(conversation, bulles de messages avec rendu Markdown, cartes de sources, visionneuse de
figures, sélecteur de périmètre documentaire), \texttt{components/layout} et
\texttt{components/ui} (coquille applicative, navigation, composants réutilisables, bascule
de thème clair/sombre). Un client API central (\texttt{api/client.ts}) joint le jeton Clerk
à chaque appel ; le hook \texttt{useChat} encapsule l'état conversationnel ; le routage est
assuré par React Router avec routes protégées.

\rapportfig[0.9]{ui_chat.png}{Interface conversationnelle : réponse sourcée avec citations
et figure affichée sous la réponse.}{fig:ui-chat}

\rapportfig[0.9]{ui_documents.png}{Vue de gestion du corpus : groupes documentaires, état
d'ingestion en temps réel et actions par document.}{fig:ui-documents}

\section{Communication et flux entre composants}
\label{sec:flux}

Tous les échanges frontend--backend sont des appels REST JSON authentifiés par jeton
Bearer, hormis les images servies statiquement (\texttt{/static/images/...}). Les deux flux
majeurs sont modélisés en séquence : le flux d'ingestion
(figure~\ref{fig:seq-ingestion}) illustre le découplage entre la requête de téléversement
(synchrone, immédiate) et le traitement (asynchrone, séquentiel) ; le flux de question
(figure~\ref{fig:seq-chat}) montre l'enchaînement routage automatique -- recherche hybride
-- génération -- restitution sourcée.

\rapportfig[0.98]{uml_sequence_ingestion.pdf}{Diagramme de séquence du flux d'ingestion
d'un document.}{fig:seq-ingestion}

\rapportfig[0.98]{uml_sequence_chat.pdf}{Diagramme de séquence du flux de
question--réponse sourcée.}{fig:seq-chat}

\section{Modélisation UML}
\label{sec:uml}

Outre les diagrammes de séquence ci-dessus, deux vues complètent la modélisation. Le
diagramme de cas d'utilisation (figure~\ref{fig:usecase}) recense les interactions de
l'utilisateur authentifié avec le système ; le diagramme de classes du modèle de données
(figure~\ref{fig:modele-donnees}) est présenté en section suivante.

\rapportfig[0.85]{uml_cas_utilisation.pdf}{Diagramme de cas d'utilisation de l'assistant
RAG multimodal.}{fig:usecase}

\section{Modèle de données}
\label{sec:modele-donnees}

La persistance repose sur une séparation stricte entre \emph{métadonnées relationnelles}
(SQLite) et \emph{contenu vectorisé} (ChromaDB), reliés par l'identifiant de document
(UUID) :

\begin{itemize}
  \item \textbf{SQLite}, quatre tables : \texttt{documents} (identité, groupe, état
        d'ingestion, étape courante, message d'erreur, chemin du fichier, lien
        \texttt{source\_document\_id} pour les duplicatas), \texttt{groups} (groupes
        documentaires nommés, avec groupe par défaut non supprimable),
        \texttt{conversations} (rattachées à un \texttt{user\_id} Clerk) et
        \texttt{messages} (rôle, contenu, sources sérialisées en JSON, horodatage) ;
  \item \textbf{ChromaDB}, une collection unique \texttt{documents} où chaque entrée
        porte : le vecteur (1024 dim.), le texte (chunk ou description de figure) et les
        métadonnées \texttt{\{document\_id, document\_name, page, chunk\_type
        (text|image), image\_path\}} ; c'est ce champ \texttt{chunk\_type} qui matérialise
        l'indexation multimodale unifiée ;
  \item \textbf{Système de fichiers} : \texttt{uploads/} (originaux),
        \texttt{markdown/} (exports), \texttt{images/<document\_id>/} (figures extraites).
\end{itemize}

Un point de conception notable : la \emph{duplication} d'un document ne copie ni fichier,
ni vecteurs, ni images : le duplicata référence le document racine via
\texttt{source\_document\_id}, et la consultation comme la suppression résolvent ce lien
(seule la suppression du document racine détruit le contenu partagé).

\rapportfig[0.95]{uml_modele_donnees.pdf}{Modèle de données : tables SQLite, collection
ChromaDB et stockage fichiers, reliés par l'identifiant de document.}{fig:modele-donnees}

\section{Justification des choix technologiques}
\label{sec:choix-technos}

\begin{table}[!htb]
\centering\small
\caption{Choix technologiques et justifications.}
\label{tab:choix}
\begin{tabularx}{\textwidth}{@{}p{3.1cm}p{3.4cm}X@{}}
\toprule
\textbf{Besoin} & \textbf{Choix} & \textbf{Justification} \\
\midrule
Analyse documentaire & Docling & Open source, détection de mise en page et TableFormer par apprentissage profond, extraction native des figures avec légendes, chunker sémantique intégré ; modèles téléchargeables hors ligne (indispensable derrière le proxy) \\
Backend PDF & pypdfium2 & Le backend Docling par défaut provoquait des \texttt{std::bad\_alloc} sur des PDF à images volumineuses ; pypdfium2 est plus robuste en mémoire \\
Embeddings & \texttt{mistral-embed} & Seule API d'embeddings de qualité \emph{accessible à travers Zscaler} (Hugging Face bloqué) ; 1024 dim., multilingue, batch natif ; backends alternatifs (HF Inference, local transformers) conservés et commutables par configuration \\
Captioning VLM & Gemma-3 27B (DeepInfra) & Excellent rapport qualité/coût pour la description de figures techniques ; accessible via l'API unifiée Hugging Face Inference Providers \\
Base vectorielle & ChromaDB & Embarquée et persistante : zéro serveur à installer (PC verrouillé), portage OCI trivial ; API de filtrage par métadonnées suffisante au besoin \\
Retrieval lexical & \texttt{rank\_bm25} + RRF & Bibliothèque pure Python ; RRF fusionne les classements sans calibration de scores (constante $k=60$ de Cormack et al.) \\
LLM de génération & \texttt{mistral-small-latest} & Qualité suffisante pour la synthèse contrainte, latence et coût faibles, API stable et accessible \\
API backend & FastAPI & Validation Pydantic, documentation OpenAPI automatique, montage de fichiers statiques, écosystème Python (Docling, ChromaDB) \\
Frontend & React 18 + Vite + TypeScript & SPA réactive, typage de bout en bout avec les schémas Pydantic, rechargement instantané \\
Authentification & Clerk & Authentification déléguée complète (UI incluse) sans stockage de mots de passe ; vérification backend purement cryptographique (JWKS), sans appel réseau par requête \\
Cache & Redis / FakeRedis & Même code en dev (FakeRedis en mémoire, aucun serveur) et en prod (OCI Cache), commutation par variable d'environnement ; conception \emph{fail-open} \\
Évaluation & RAGAS + Groq (Llama 3.3 70B) & Métriques standardisées ; juge indépendant du système évalué pour éviter le biais d'auto-évaluation \\
\bottomrule
\end{tabularx}
\end{table}

Deux arbitrages méritent un développement :

\paragraph{Sans LangChain.} Le pipeline est assemblé directement sur les SDK natifs
(\texttt{mistralai}, \texttt{huggingface\_hub}, \texttt{chromadb}). Pour un pipeline de
cette taille, les abstractions de LangChain auraient ajouté une épaisseur difficile à
auditer sans rien apporter : chaque étape (chunking, embeddings, fusion, prompt) est ici
un module Python explicite de quelques dizaines de lignes, directement instrumentable,
un atout décisif lors du diagnostic de qualité du retrieval.

\paragraph{Captioning plutôt qu'embeddings multimodaux.} L'approche par description VLM
(section~\ref{sec:approches}) a été préférée à CLIP/ColPali : elle conserve un index
unique et un pipeline de retrieval inchangé, produit des descriptions réutilisées telles
quelles dans l'interface (galerie de figures, cartes de sources), et ne requiert aucun
modèle auto-hébergé, rédhibitoire sur le poste de développement.

\section{Conclusion}

L'architecture conçue répond aux besoins formalisés par un système à trois niveaux, sobre
et entièrement auditable : une chaîne d'ingestion asynchrone sérialisée qui transforme des
documents hétérogènes en un index multimodal unifié, et une chaîne de restitution qui
combine routage LLM, recherche hybride RRF et génération contrainte citée. Chaque choix
technologique a été arbitré au regard du triptyque qualité / contraintes d'environnement /
portabilité OCI. Le chapitre suivant décrit la mise en œuvre concrète de cette
architecture et sa validation expérimentale.
